\subsection{The existing paper projective model and the actual native scheme}
\label{subsec:native-projective-paper-bridge}

\begin{theorem}[The paper model is the trivial-base normalized spectrum]
\label{projective:paper-normalized-spectrum-bridge}
\lean{EnvelopeProperty.ProjBerk.normalizedSpectrumHomeomorph,
  EnvelopeProperty.ProjBerk.normalizedSpectrumHomeomorph_val,
  EnvelopeProperty.ProjBerk.normalizedSpectrumHomeomorph_val_apply,
  EnvelopeProperty.ProjBerk.normalizedSpectrumHomeomorph_mem_chartOpen,
  EnvelopeProperty.ProjBerk.normalizedSpectrumChartHomeomorph,
  EnvelopeProperty.ProjBerk.normalizedSpectrumHomeomorph_coordinateChartMap}
\leanok
\uses{spectrum:trivial-field-bridge,projective:general-relative-normalized-charts}
For a trivially valued field $k$, the existing normalized graded model
$\operatorname{ProjBerk}(k,n)$ is canonically homeomorphic to the general
normalized projective spectrum above the trivial seminorm on $k$.
Every polynomial value, every coordinate-open condition and the entire
affine ratio-chart seminorm are preserved.
\end{theorem}
\begin{proof}
The genuine affine homeomorphism between the existing affine
analytification and the relative spectrum above the trivial field
seminorm preserves the entire polynomial valuation. Restrict it to the
actual gradedness and coordinate-normalization predicates, which depend
on these same polynomial values. Reorder the fixed coefficient and
normalization subtypes using the proved coefficient homeomorphism.
Both maps retain the original valuation, so all polynomial evaluations
and coordinate-open conditions agree. Restrict to each such open.
The two whole ratio-chart maps are the pullbacks of the same genuine
localized valuation under the same actual affine ratio homomorphism;
localization extension and proof irrelevance identify them on every
polynomial, proving the entire chart square.
\end{proof}

\begin{theorem}[The paper model has the genuine native projective interpretation]
\label{projective:paper-native-projective-homeomorphism}
\lean{EnvelopeProperty.ProjBerk.nativeHomeomorph,
  EnvelopeProperty.ProjBerk.nativeHomeomorph_apply,
  EnvelopeProperty.ProjBerk.nativeHomeomorph_support_toIdeal,
  EnvelopeProperty.ProjBerk.nativeHomeomorph_point,
  EnvelopeProperty.ProjBerk.nativeHomeomorph_chart,
  EnvelopeProperty.ProjBerk.nativeHomeomorph_val_apply,
  EnvelopeProperty.ProjBerk.nativeHomeomorph_eval_form_of_val_X_eq_one,
  EnvelopeProperty.ProjBerk.nativeHomeomorph_log_form_of_val_X_eq_one}
\leanok
\uses{projective:paper-normalized-spectrum-bridge,projective:normalized-native-homeomorphism}
The existing paper model is canonically homeomorphic to the actual native
valued-point space of the genuine polynomial Proj scheme over the whole
trivially valued base point. Its actual homogeneous prime support remains
the full original polynomial seminorm kernel. Every actual coordinate
polynomial section has the original affine chart value. For a homogeneous
form $F$ of degree $d$, a coordinate of value one gives its native section
$F/X_j^d$ the original normalized value of $F$. Its actual logarithm
equals the existing paper function
$\log|F|$, including zero and negative infinity.
\end{theorem}
\begin{proof}
Compose the preceding genuine whole-valuation homeomorphism with the
proved canonical native homeomorphism for the trivial coefficient
seminorm. The canonical native point uses the actual degree-zero
projective stalk and its actual residue field, so its support is the
original homogeneous kernel. The full generic coordinate-chart square
and the preceding whole affine specialization identify all actual
polynomial section evaluations. Apply the genuine homogeneous-fraction
value formula when the chosen coordinate has value one; its denominator
has value one, giving the claimed form value. Applying the genuine
extended logarithm proves equality with the existing logarithmic form
function, including zero-valued forms.

This proves the point-space and topology interpretation used by the
existing paper projective interfaces. It preserves their field API. The
actual projective residue field remains distinct from the generally
larger affine-cone residue field. An analytic structure sheaf and the
identification of coherent sheaf model ideals are separate constructions.
\end{proof}
